\chapter{Estado da Arte}
\label{cap:estado_da_arte}

% TODO: Introdução curta — o objetivo deste capítulo é enquadrar as decisões técnicas
% do estágio no contexto das práticas e ferramentas existentes, com foco em arquitetura
% de backend, qualidade de software e entrega contínua.

%----------------------------------------------------------------------------------------
\section{Trabalhos Relacionados}
\label{sec:trabalhos_relacionados}

% Explorar dois eixos principais de trabalho relacionado:
%
% 1. PLATAFORMAS DE ASSISTÊNCIA DIGITAL A IDOSOS
%    Analisar brevemente soluções existentes como ElliQ, Care.Coach e plataformas de
%    telemedicina em contexto de ERPI (Estruturas Residenciais para Pessoas Idosas).
%    O objetivo não é um estudo exaustivo destas plataformas, mas sim identificar
%    a lacuna que o CompyAI preenche: a ausência de um canal de comunicação natural
%    (voz) combinado com geração automática de informação estruturada para cuidadores.
%    Esta análise justifica a existência do produto, não as escolhas técnicas.
%
% 2. ARQUITETURAS DE BACKEND PARA SISTEMAS MULTI-TENANT
%    Rever trabalhos e práticas estabelecidas sobre arquitetura de APIs REST em sistemas
%    com múltiplos perfis de utilizador e isolamento de dados por instituição (multi-tenant).
%    Referir padrões como Clean Architecture \parencite{martin2017clean} e a separação
%    entre camada de apresentação, lógica de negócio e acesso a dados.
%    Relacionar com o contexto do CompyAI: três perfis de utilizador (admin, caregiver,
%    senior) com permissões distintas e dados partilhados por instituição.
%
% Sugestão: tabela resumo comparando as soluções de assistência digital existentes
% em termos de canal de interação, integração com cuidadores e modelo de dados.

%----------------------------------------------------------------------------------------
\section{Tecnologias e Práticas Existentes}
\label{sec:tecnologias_existentes}

% Organizar em torno das decisões técnicas mais relevantes para o trabalho do estágio.
% Para cada área, comparar alternativas e justificar a escolha adotada.
%
% 1. FRAMEWORKS DE BACKEND
%    Comparar NestJS, Express e Fastify em termos de estrutura, injeção de dependências,
%    suporte a TypeScript e ecossistema de módulos.
%    Justificar a escolha do NestJS pela estrutura modular que favorece separação de
%    responsabilidades, pelo sistema de injeção de dependências nativo (que facilita
%    testes unitários) e pela maturidade do ecossistema para APIs REST em TypeScript.
%
% 2. ORM E ACESSO A DADOS
%    Comparar Drizzle ORM, Prisma e TypeORM em termos de type-safety, performance de
%    queries, controlo sobre as migrações e compatibilidade com PostgreSQL.
%    Justificar a escolha do Drizzle pela abordagem SQL-first (o schema é definido em
%    TypeScript mas as queries ficam próximas do SQL) e pelo controlo granular sobre
%    migrações — relevante para um ambiente de produção onde as migrações são aplicadas
%    automaticamente no arranque do contentor.
%
% 3. AUTENTICAÇÃO E AUTORIZAÇÃO
%    Comparar abordagens de autenticação stateless (JWT) vs. stateful (sessões em base
%    de dados) para APIs REST com múltiplos perfis de utilizador.
%    Justificar a escolha de JWT pela adequação a uma arquitetura de serviços
%    independentes (o token carrega o perfil do utilizador, eliminando consultas à base
%    de dados em cada pedido) e pela facilidade de integração com guards do NestJS para
%    implementar RBAC (Role-Based Access Control).
%
% 4. ESTRATÉGIAS DE TESTES EM BACKEND NODE.JS
%    Comparar Jest e Vitest como frameworks de testes para projetos TypeScript/ESM.
%    Justificar a escolha do Vitest pela compatibilidade nativa com ESM, pela velocidade
%    de execução e pela API compatível com Jest (curva de aprendizagem reduzida).
%    Referir a pirâmide de testes \parencite{cohn2009succeeding} como modelo orientador
%    da estratégia adotada.
%
% 5. CONTENTORIZAÇÃO E ENTREGA CONTÍNUA
%    Comparar abordagens de deployment: servidor dedicado sem contentores, Docker com
%    Compose, e orquestradores como Kubernetes.
%    Justificar a escolha de Docker Compose para uma startup em fase inicial: complexidade
%    operacional reduzida face ao Kubernetes, infraestrutura declarativa e versionada,
%    e facilidade de recriar o ambiente de produção localmente.
%    Referir \textcite{humble2010continuous} sobre os princípios de entrega contínua que
%    orientaram a automação das migrações e o processo de atualização da plataforma.
%
% 6. INTELIGÊNCIA ARTIFICIAL (menção breve)
%    Contextualizar brevemente o papel dos fornecedores externos de IA (STT, LLM, TTS)
%    na plataforma, sem aprofundar as comparações tecnológicas — o foco do estágio foi
%    a integração segura destes serviços no backend, não a sua avaliação comparativa.
